\documentclass[a4paper,11pt]{article}
\usepackage[T1]{fontenc}
\usepackage[polish]{babel}
\usepackage{lmodern,amsmath,booktabs}
\usepackage[margin=2.6cm]{geometry}
\title{Okres wahadła}
\author{Przykład dydaktyczny}
\date{}
\begin{document}
\maketitle
\section{Model}
Dla małych wychyleń okres wahadła o długości $L$ opisuje wzór:
\[
T=2\pi\sqrt{\frac{L}{g}}.
\]
Przyjmujemy $g=9{,}81\,\mathrm{m/s^2}$.
\section{Dane przykładowe}
\begin{itemize}
\item Wydłużenie nici zwiększa okres.
\item Dane służą do ćwiczenia; nie są wynikiem pomiaru.
\end{itemize}
\begin{center}
\begin{tabular}{cc}
\toprule
Długość $L$ (m) & Okres $T$ (s) \\
\midrule
$0{,}20$ & $0{,}90$ \\
$0{,}40$ & $1{,}27$ \\
\bottomrule
\end{tabular}
\end{center}
Wartości w tabeli można zestawić z modelem, ale dwa wiersze nie
wystarczą do oceny jego dokładności.
\end{document}
